\documentclass[exercice]{thedocument}%
\usepackage{texmaths-tikz}
\startdatakeys
    \source{inconnue}
    \BO{2020}
    \cycle{4}
    \competencesDuSocle{RE}
    \competencesTravaillees{A1a2, D2c4, D2a*}
\enddatakeys
\begin{document}%
    Construire l'image du triangle SAM par l'homothétie de centre O et de rapport
    $\dfrac{1}{4}$ :

    \vspace{0.4em}
    \begin{center}
    \begin{tikzpicture}[x=0.5cm, y=0.5cm]
        % Grille
        \draw[gridcol, very thin, step=1] (-5,-2) grid (9,5);

        % Point O  (-4, 0)
        \ptmark{-4}{0}
        \node[ptcol, above left, font=\small\itshape] at (-4,0) {\mathspoint{O}};

        % Point S  (0, 4)
        \ptmark{0}{4}
        \node[ptcol, above right, font=\small\itshape] at (0,4) {\mathspoint{S}};

        % Point M  (0, 0)
        \ptmark{0}{0}
        \node[ptcol, below right, font=\small\itshape] at (0,0) {\mathspoint{M}};

        % Point A  (8, 4)
        \ptmark{8}{4}
        \node[ptcol, above right, font=\small\itshape] at (8,4) {\mathspoint{A}};

        % Triangle SAM  (angle droit en S)
        \draw[black!80, line width=1.5pt] (0,4) -- (8,4) -- (0,0) -- cycle;
    \end{tikzpicture}
    \end{center}
\end{document}